% GENERATED FILE. Edit canonical lessons, then run npm run generate:books.
% canonical-source-hash: fnv1a64:0a2e87d7ebbb2f05
% canonical-lessons: IT-C342-ammettere, IT-C342-concordare, IT-C342-consenso, IT-C342-compromesso, IT-C342-cedere

\chapter{Trovare un accordo}
\label{ch:it-trovare-un-accordo}
\hlchaptermodality{342}

\begin{chapteropening}
\textbf{By the end of this chapter:} \emph{I can admit someone is right, agree with them, give way and reach a compromise.}

Complete the last lesson of chapter 342: I can admit someone is right, agree with them, give way and reach a compromise.
\end{chapteropening}

\section[ammettere]{ammettere — to admit}

\label{lesson:IT-C342-ammettere}



\begin{quote}
Before the new one: say the Italian for to answer back, then the Italian for to contradict.
\end{quote}

\subsection*{You'll want to know: ammettere}
\textbf{ammettere} — \textquotedblleft{}to admit\textquotedblright{}.

\begin{quote}
Devo ammettere che il tuo piano era migliore.
\end{quote}

\begin{quote}
\emph{I have to admit that your plan was better.}
\end{quote}

Conjugated like \emph{mettere}; past participle \emph{ammesso}. It also means to let someone in: \emph{ammettere a un corso}.

\begin{cousinweb}
\emph{ammettere} is inherited from Latin \emph{admittere} \textquotedblleft{}to let in, to allow\textquotedblright{} (EtymDB), built from \emph{ad} \textquotedblleft{}to\textquotedblright{} and \emph{mittere} \textquotedblleft{}to send\textquotedblright{}: you let a fact in.
\end{cousinweb}

\begin{grammarlens}[title={What you've built}]
One more word for telling, explaining and describing at B1.
\end{grammarlens}

\subsection*{Your turn}
\begin{itemize}
\raggedright
  \item \emph{Say it:} \emph{ammettere}
  \item \emph{Say it:} \emph{ammettere}, once more
  \item \emph{Say it:} a sentence of your own with \emph{ammettere}
  \item \emph{Recall:} say the Italian for to answer back, then the Italian for to contradict, then say \emph{ammettere} again
\end{itemize}

\begin{tcolorbox}[breakable,colback=teal!4,colframe=teal!35!black,title={Before you move on}]
What does ammettere mean? (\textquotedblleft{}To admit\textquotedblright{}.) Say it once more.
\end{tcolorbox}
\section[concordare]{concordare — to agree (on), to settle together}

\label{lesson:IT-C342-concordare}



\begin{quote}
Before the new one: say the Italian for to contradict, then the Italian for to admit.
\end{quote}

\subsection*{You'll want to know: concordare}
\textbf{concordare} — \textquotedblleft{}to agree (on), to settle together\textquotedblright{}.

\begin{quote}
Abbiamo concordato il prezzo con il proprietario.
\end{quote}

\begin{quote}
\emph{We agreed on the price with the owner.}
\end{quote}

\emph{Concordare con qualcuno} — to agree with someone; \emph{concordare una data} — to agree on a date. A little more formal than \emph{essere d'accordo}.

\begin{cousinweb}
\emph{concordare} is a learned borrowing from Latin \emph{concordāre} \textquotedblleft{}to agree, to be of one mind\textquotedblright{} (Lewis \& Short), from \emph{concors} \textquotedblleft{}of one heart\textquotedblright{}, built from \emph{con-} \textquotedblleft{}together\textquotedblright{} and \emph{cor, cordis} \textquotedblleft{}heart\textquotedblright{}.
\end{cousinweb}

\begin{grammarlens}[title={What you've built}]
One more word for telling, explaining and describing at B1.
\end{grammarlens}

\subsection*{Your turn}
\begin{itemize}
\raggedright
  \item \emph{Say it:} \emph{concordare}
  \item \emph{Say it:} \emph{concordare}, once more
  \item \emph{Say it:} a sentence of your own with \emph{concordare}
  \item \emph{Recall:} say the Italian for to contradict, then the Italian for to admit, then say \emph{concordare} again
\end{itemize}

\begin{tcolorbox}[breakable,colback=teal!4,colframe=teal!35!black,title={Before you move on}]
What does concordare mean? (\textquotedblleft{}To agree (on), to settle together\textquotedblright{}.) Say it once more.
\end{tcolorbox}
\section[il consenso]{il consenso — the consent; the agreement (of many)}

\label{lesson:IT-C342-consenso}



\begin{quote}
Before the new one: say the Italian for to admit, then the Italian for to agree.
\end{quote}

\subsection*{You'll want to know: il consenso}
\textbf{il consenso} — \textquotedblleft{}the consent; the agreement (of many)\textquotedblright{}.

\begin{quote}
Non puoi usare la mia foto senza il mio consenso.
\end{quote}

\begin{quote}
\emph{You cannot use my photo without my consent.}
\end{quote}

\emph{Dare il proprio consenso} — to give your consent. It can also be broad agreement in a group.

\begin{cousinweb}
\emph{consenso} is a learned borrowing from Latin \emph{cōnsēnsus} \textquotedblleft{}agreement, accord\textquotedblright{} (Lewis \& Short), from \emph{cōnsentīre} \textquotedblleft{}to agree\textquotedblright{}, built from \emph{con-} \textquotedblleft{}together\textquotedblright{} and \emph{sentīre} \textquotedblleft{}to feel\textquotedblright{}.
\end{cousinweb}

\begin{grammarlens}[title={What you've built}]
One more word for telling, explaining and describing at B1.
\end{grammarlens}

\subsection*{Your turn}
\begin{itemize}
\raggedright
  \item \emph{Say it:} \emph{il consenso}
  \item \emph{Say it:} \emph{il consenso}, once more
  \item \emph{Say it:} a sentence of your own with \emph{il consenso}
  \item \emph{Recall:} say the Italian for to admit, then the Italian for to agree, then say \emph{il consenso} again
\end{itemize}

\begin{tcolorbox}[breakable,colback=teal!4,colframe=teal!35!black,title={Before you move on}]
What does il consenso mean? (\textquotedblleft{}The consent; the agreement (of many)\textquotedblright{}.) Say it once more.
\end{tcolorbox}
\section[il compromesso]{il compromesso — the compromise}

\label{lesson:IT-C342-compromesso}



\begin{quote}
Before the new one: say the Italian for to agree, then the Italian for the consent.
\end{quote}

\subsection*{You'll want to know: il compromesso}
\textbf{il compromesso} — \textquotedblleft{}the compromise\textquotedblright{}.

\begin{quote}
Alla fine abbiamo trovato un compromesso: un film e poi una pizza.
\end{quote}

\begin{quote}
\emph{In the end we found a compromise: a film and then a pizza.}
\end{quote}

\emph{Arrivare a un compromesso} — to reach a compromise. It is also the preliminary contract when you buy a house.

\begin{cousinweb}
\emph{compromesso} comes from Latin \emph{compromissum} \textquotedblleft{}a mutual promise to abide by an arbiter's decision\textquotedblright{} (Lewis \& Short), from \emph{compromittere}, built on \emph{prōmittere} \textquotedblleft{}to promise\textquotedblright{} and \emph{mittere} \textquotedblleft{}to send\textquotedblright{}.
\end{cousinweb}

\begin{grammarlens}[title={What you've built}]
One more word for telling, explaining and describing at B1.
\end{grammarlens}

\subsection*{Your turn}
\begin{itemize}
\raggedright
  \item \emph{Say it:} \emph{il compromesso}
  \item \emph{Say it:} \emph{il compromesso}, once more
  \item \emph{Say it:} a sentence of your own with \emph{il compromesso}
  \item \emph{Recall:} say the Italian for to agree, then the Italian for the consent, then say \emph{il compromesso} again
\end{itemize}

\begin{tcolorbox}[breakable,colback=teal!4,colframe=teal!35!black,title={Before you move on}]
What does il compromesso mean? (\textquotedblleft{}The compromise\textquotedblright{}.) Say it once more.
\end{tcolorbox}
\section[cedere]{cedere — to give in, to give way}

\label{lesson:IT-C342-cedere}



\begin{quote}
Before the new one: say the Italian for the consent, then the Italian for the compromise.
\end{quote}

\subsection*{You'll want to know: cedere}
\textbf{cedere} — \textquotedblleft{}to give in, to give way\textquotedblright{}.

\begin{quote}
Dopo un'ora di discussione, mia sorella ha ceduto.
\end{quote}

\begin{quote}
\emph{After an hour of discussion, my sister gave in.}
\end{quote}

\emph{Cedere il posto} — to give up your seat. \emph{Non cedere} — to hold firm.

\begin{cousinweb}
\emph{cedere} is inherited from Latin \emph{cēdere} \textquotedblleft{}to go, to withdraw, to yield\textquotedblright{} (EtymDB): you step back and let the other side pass.
\end{cousinweb}

\begin{grammarlens}[title={What you've built}]
One more word for telling, explaining and describing at B1.
\end{grammarlens}

\subsection*{Your turn}
\begin{itemize}
\raggedright
  \item \emph{Say it:} \emph{cedere}
  \item \emph{Say it:} \emph{cedere}, once more
  \item \emph{Say it:} a sentence of your own with \emph{cedere}
  \item \emph{Recall:} say the Italian for the consent, then the Italian for the compromise, then say \emph{cedere} again
\end{itemize}

\begin{tcolorbox}[breakable,colback=teal!4,colframe=teal!35!black,title={Before you move on}]
What does cedere mean? (\textquotedblleft{}To give in, to give way\textquotedblright{}.) Say it once more.
\end{tcolorbox}
